
\documentclass[english]{article}
%\documentclass[aps,pra,twocolumn,english]{revtex4}
%\documentclass[english]{iopart}
\usepackage[T1]{fontenc}
\usepackage[latin9]{inputenc}
\usepackage{babel}
\usepackage{graphicx,amsmath,amssymb,color}
\usepackage[normalem]{ulem}
\usepackage{amsfonts}
\usepackage[toc,page]{appendix}
%\usepackage[ocgcolorlinks,colorlinks=true,linkcolor=blue,citecolor=red]{hyperref}
\usepackage{hyperref}
\usepackage{latexsym}
\usepackage{amsfonts}
\usepackage{algpseudocode}
\usepackage{amsthm}
\usepackage{mathrsfs}
\usepackage{natbib}
\usepackage{color,verbatim}
\usepackage{psfrag}
\bibliographystyle{unsrt}
%\bibliographystyle{acm}
%\bibliographystyle{unsrtnat}
%\usepackage[style=numeric]{biblatex}

\usepackage{subcaption} % allows for subfigures


\usepackage[svgnames]{xcolor}
\usepackage{tikz}

\usepackage{algorithm}


%\input{tikzit/tikz_preamble.tex}  % load separately defined tikz preamble


\newcommand{\bra}[1]{\langle#1 |}
\newcommand{\ket}[1]{|#1 \rangle}
\newcommand{\braket}[2]{\left \langle #1 \mid #2 \right \rangle}
\newcommand{\bigbraket}[2]{\left \langle #1 \middle\vert #2 \right \rangle}
\newcommand{\ketbra}[2]{\vert #1 \rangle \! \langle #2 \vert}
\newcommand{\overlap}[2]{\left \langle #1 | #2 \right\rangle}
\newcommand{\average}[1]{\langle #1 \rangle}
\newcommand{\sandwich}[3]{\left \langle #1 \mid #2 \mid #3 \right\rangle}
\newcommand{\innerprod}[2]{\left \langle #1 | #2 \right\rangle}

\begin{document}


\title{Wavepackets and multi (frequency) mode states}
\author{Nicholas Chancellor}


\maketitle

\textcolor{red}{Assumed Raouf has already done quantisation of fields if not it needs to come before this}

\today % added for drafting, remove in final version


{\small Disclaimer: lecture notes may contain ``typo'' type errors (factors of $2$, minus signs, etc...) if something like this looks wrong there is a good chance it is, please let me know}

\section{Wavepackets of light}

\begin{itemize}
\item Electric field operator for a single mode (G+K eq.~2.127): 
\begin{equation}
\hat{\vec{E}}(z)=i\left(\frac{\hbar \omega}{2 \epsilon_0 V}\right)^\frac{1}{2}\vec{e}_x[\hat{a}e^{i\vec{k}\cdot \vec{r}-i \omega t}-\hat{a}^\dagger e^{-i\vec{k}\cdot\vec{r}+i \omega t}]
\end{equation}
\item Only think about direction of propagation and a single polarisation, 1D problem
\begin{equation}
\hat{E}(z)=i\left(\frac{\hbar \omega}{2 \epsilon_0 V}\right)^\frac{1}{2}[\hat{a}e^{i k z-i \omega t}-\hat{a}^\dagger e^{-ik z+i \omega t}]
\end{equation}
\item Infinite physical extent (in three dimensions, but let's focus on the direction of propagation)... but our devices have finite size, we need more modes!
\item Multimode field in 1D 
\begin{equation}
\hat{E}(z)=i\int_kdk\left(\frac{\hbar \omega}{2 \epsilon_0 V}\right)^\frac{1}{2}[\hat{a}_ke^{i k z-i \omega_k t}-\hat{a}_k^\dagger e^{-ik z+i \omega_k t}]
\end{equation}
\item How do we build a wavepacket of finite extent $\rightarrow$ if we take $\average{E}$ (and also note that in free space $\omega_k=c k$ where $c$ is the speed of light)
\begin{equation}
\average{\hat{E}(z)}=i\int_kdk\left(\frac{\hbar \omega}{2 \epsilon_0 V}\right)^\frac{1}{2}[\average{\hat{a}_k}e^{i k z-i c k t}-\average{\hat{a}_k^\dagger} e^{-ik z+i c k t}]
\end{equation}
\begin{itemize}
\item Propagates in the $+z$ direction at the speed of light $c$ by definition, replace $z$ with $z(t)=z-ct$
\begin{equation}
\average{\hat{E}(z)}=i\int_kdk\left(\frac{\hbar \omega}{2 \epsilon_0 V}\right)^\frac{1}{2}[\average{\hat{a}_k}e^{i k z(t)}-\average{\hat{a}_k^\dagger} e^{-ik z(t)}]
\end{equation}
\item Fun side-note/exercise, if we instead had $\omega_k= \frac{\hbar k^2}{2 m}$, this would describe a (non-relativistic) massive particle, and speed would be $\frac{\partial \omega}{\partial k}=\frac{\hbar k}{m}$, the relationship between $\omega$ and $k$ is known as the \textbf{dispersion relation} and is fundamental in almost all wave physics, light can also have more complicated dispersion in some materials
\end{itemize}
\item Then this looks like a Fourier transform with $\average{\hat{a}_k}$ and $\average{\hat{a}_k^\dagger}$ as the Fourier coefficients $\rightarrow$ we can make wavepackets of arbitary shape; 
Problem? $E$ is a physical observable, it must be real valued everywhere
\begin{itemize}
\item $\average{\hat{a}_k}$ and $\average{\hat{a}_k^\dagger}$ are not independent!
\item Consider $\ket{\psi}=C_n\ket{n}+C_{n+1}\ket{n+1}$, $\average{\hat{a}_k}_{\psi}=C_nC^\star_{n+1}$ but $\average{\hat{a}_k^\dagger}_\psi=C^\star_nC_{n+1}$
\item Follows that $\average{\hat{a}_k^\dagger}=\average{\hat{a}_k}^\star$
\item Using the definition of sin and cosine\footnote{In case you have forgotten, or not seen, these definitions: $\sin(x)=\frac{\exp(ix)-\exp(-ix)}{2i}$ and $\cos(x)=\frac{\exp(ix)+\exp(-ix)}{2}$ } (and the fact that cosine is an even function): 
\begin{equation}
\average{\hat{E}(z)}=2\int_kdk\left(\frac{\hbar \omega}{2 \epsilon_0 V}\right)^\frac{1}{2}[\mathrm{Re}(\average{\hat{a}_k}) \sin(k z-  c k t)-\mathrm{Im}(\average{\hat{a}_k}) \cos(k z-  c k t)]
\end{equation}
\item Purely real by definition
\end{itemize}
\end{itemize}
\section{Multimode wavepackets in practice}
\begin{itemize}
\item A problem and its resolution:
\begin{itemize}
\item Potential problem? Don't operators describing different frequency modes commute $[\hat{a}_k,\hat{a}_{k'}]=0$ for $k\neq k'$? since we have a continuum of $k$ how do we see any quantum effects at all?
\item Resolution: Instead of treating creation/annihilation operators for single modes, consider creation/annihilation for multimode wavepackets $[\hat{a}[\Psi],\hat{a}[\Phi]]\propto \braket{\psi}{\phi}$, where $\hat{a}[\Psi]=\int dk\Psi(k) a_k$, overlap of wavefunctions is \textbf{distinguishability} $\rightarrow$ identical wavepackets (including phases, which indicate when they arrive) are indistinguishable
\item Making indistinguishable wavepackets is key to many quantum optics applications, such as Boson sampling
\item ``Easy''\footnote{relative to other experiments, none of this is really easy} to measure indistinguishability technically  with single photon wave packets, send both wavepackets into a beam splitter, measure how many times detectors click simultaneously, Hong-Ou-Mandel effect (see: \href{arXiv:quant-ph/0212017}{https://arxiv.org/abs/quant-ph/0212017})
\end{itemize}
 An uncertainty relationship between $k$ and $z$:
\item Taking the Fourier transform of $k$ gives $k=-i\frac{\partial}{\partial z}$, taking $\average{\Delta k}=\average{k^2}-\average{k}^2=-\int_z dz\frac{\partial^2E^2(z)}{\partial z^2}+\int_zdz(\frac{\partial E^2(z)}{\partial z})^2$ 
\begin{itemize}
\item n.b.~ I have used $\average{E^2(z)}$ in my definition so it even works for states like number states where $\average{E(z)}$ is zero everywhere.\footnote{This has the complication that $\average{E^2(z)}$ is non-zero for vacuum, but this can be mitigated by integrating over a finite extent and considering states of many photons so the vacuum contribution is small}
\end{itemize}
\item Conceptually: a narrow peak requires a large negative curvature, therefore smaller $\Delta z$ leads to larger $\Delta k$
\item Mathematically: $[k,z]=-i\frac{\partial}{\partial  z} z+i z \frac{\partial}{\partial z}$ by the product rule $-i\frac{\partial}{\partial  z}  (z E^2(z))=-iE^2(z)-i z \frac{\partial}{\partial z}$, cancelling terms and taking expectation $\average{[k,z]}=-i\int_zdzE^2(z)$
\item Therefore $\average{\Delta k \Delta z}\ge \int_zdzE^2(z)$ (note that $E^2(z)$ is not normalized like a wavefunction so can't just cancel to $1$)
\item Converse is not necessarily true, possible to have large (or even infinite) $\Delta k$ and $\Delta z$ simultaneously
\item But wait! none of this depended on the system being quantum, isn't this just a statement about classical Fourier transforms which would also hold for very classical systems like water waves? Yes and no
\begin{itemize}
\item The relationship between spread of position and momentum holds for classical systems
\item Can be thought of as a ``classical uncertainty principle'' but arguably not a ``real'' uncertainty principle since classical systems do not have projective measurements so all quantities can be known simultaneously
\end{itemize}
\item Visible light has wavelength ($\lambda$) of the order 100s of nanometers $\sim 10^{-7} m$, plenty of space in experiments for many wavepackets even if $\Delta z \gg \lambda$ 
\item Lots of other fun Fourier transform identities relating $k$ and $z$ I don't think they are relevant but we can discuss if you want
\end{itemize}
\section{Photons in wavepackets}
\begin{itemize}
\item We can make a wavepacket with a single photon (a Gaussian wavepacket in this example)
\begin{equation}
\ket{\psi_\mathrm{Gauss}}=\frac{1}{\mathcal{N}}\int_kdk\exp\left[-\frac{(k-k_0)^2}{\sigma^2}\right]\hat{a}^\dagger_k\ket{0}
\end{equation}
\item Superposition of frequencies, but all states involve only $(\hat{a}^\dagger_k)^1$ so only one photon
\item Since an atom emits a photon over a finite period of time, it will always contain a superposition of frequencies (although Lorentzian/Voight, not Gaussian distributed)
\begin{equation}
\ket{\psi_\mathrm{Lorentz}}=\frac{1}{\mathcal{N}}\int_kdk \frac{\gamma^2}{\gamma^2+(k-k_0)^2}\hat{a}^\dagger_k\ket{0}
\end{equation}
\item Depending on the experiment may be justified in approximating as single frequency
\item Can also have multi-photon distributions, for example a distribution of coherent states
\begin{equation}
\ket{\psi_\mathrm{coh\,homo}}=\frac{1}{\mathcal{N}}\int_kdk \frac{\gamma^2}{\gamma^2+(k-k_0)^2} \exp(\alpha \hat{a}_k^\dagger-\alpha^\star \hat{a}_k)\ket{0}
\end{equation}
\item The above state is a bit artificial, why? Once one photon is measured all others must have the same $k$ (perfect correlation), this can happen (to a decent approximation) in a laser if all emitters are effected in the same way and is called \textbf{homogeneous broadening} 
\item If we assume photon frequencies are independent than we should have \textbf{inhomogeneous broadening} 
\begin{equation}
\ket{\psi_\mathrm{coh\,inhomo}}= \exp\left(\frac{1}{\mathcal{N}}\int_kdk\exp\left[-\frac{(k-k_0)^2}{\sigma^2}\right](\alpha \hat{a}_k^\dagger-\alpha^\star \hat{a}_k)\right)\ket{0}
\end{equation}
\item Of course we can also have both
\begin{align}
\ket{\psi_\mathrm{broad}}=\nonumber \\ \frac{1}{\mathcal{N}'}\int_{k'}dk' \frac{\gamma^2}{\gamma^2+(k'-k_0)^2} \exp\left(\int_kdk\exp\left[-\frac{(k-k')^2}{\sigma^2}\right](\alpha \hat{a}_k^\dagger-\alpha^\star \hat{a}_k)\right)\ket{0}
\end{align}
\item Which is a decent approximation of a wavepacket from a laser with Lorentzian homogenous broadening and Gaussian inhomogeneous broadening...  
\item If we wanted a continuous beam which is broadened rather than a wavepacket, we could add phases to the terms or consider an incoherent superposition of frequencies
\end{itemize}





\end{document}